\section{Discussion}
\label{sec:discussion}

\subsection{Engineering implications}

Across the assembled range the model halves the wall-superheat error of the
best classical correlation (\SI{1.84}{K} versus \SI{3.7}{K}) and is more than an
order of magnitude better on heat flux. For design this translates into a
tighter incipience margin: a thermal-margin uncertainty of $\sim\!\SI{4}{K}$
typical of correlation scatter is reduced to $\sim\!\SI{2}{K}$, allowing less
conservative operating points in microchannel electronics cooling and a
sharper boiling-incipience boundary in compact refrigerant evaporators. Because
a single model spans $\Dh = 0.35$--\SI{9.5}{mm} and $P=0.1$--\SI{16}{MPa}, the
same predictor serves applications previously requiring different correlations,
from atmospheric microchannels to high-pressure nuclear channels. As a concrete
artefact, Fig.~\ref{fig:designmap} renders the model as a design map: predicted
$\dTonb$ over the mass-flux--pressure plane for a representative stainless
channel, directly usable to place an operating point relative to incipience. The
map encodes the learned physics---superheat falling with pressure and rising
with mass flux at low pressure---in a single continuous surface.

\begin{figure}[t]
\centering
\includegraphics[width=0.72\linewidth]{fig_design_map.png}
\caption{Predicted ONB wall superheat over the mass-flux--pressure plane for a
representative stainless channel ($\Dh=\SI{3}{mm}$, $\dTsub=\SI{20}{K}$, water).
Such a map turns the model into a direct design tool for placing operating
points relative to boiling incipience.}
\label{fig:designmap}
\end{figure}

\subsection{Why an explicit reduced-pressure feature matters}

The pressure result (Section~\ref{sec:patternC}) is methodologically general.
It is tempting to assume that nondimensionalising targets and inputs with
pressure-dependent fluid properties lets a network ``see'' pressure implicitly.
We find this is not so for incipience: the property variations enter the
Reynolds, Weber and Boiling numbers too weakly for the network to recover the
first-order $P\!\uparrow\!\Rightarrow\!\dTonb\!\downarrow$ trend, and the model
remained pressure-blind until reduced pressure was supplied as a direct input.
Practitioners building data-driven boiling models should include pressure (or
reduced pressure) explicitly rather than relying on property-based
nondimensionalization to encode it.

\subsection{Physics-regularized regression versus full PINN}

The model is a physics-regularized regressor, not a field-resolving PINN. This
is a deliberate choice for sparse, heterogeneous incipience data: the scalar
ONB targets are what experiments report and what design needs, and soft Hsu and
monotonicity constraints inject the relevant physics without the data demands
and ill-posedness of full Navier--Stokes residuals on a scalar problem. Our
attempt to add a one-dimensional energy residual on an auxiliary temperature
field confirmed the trade-off: it regularised only an unsupervised field and
destabilised optimisation, with no accuracy benefit. Full field resolution
remains attractive where dense field data or trustworthy CFD are available, and
is a natural Phase~3 direction.

\subsection{Limitations}

Four limitations bound the present claims. First, the dataset, though wide in
range, is modest in size (351 points) and dominated by water; refrigerant and
non-aqueous coverage is thin, and the high-pressure end rests on a single study,
so extrapolation beyond \SI{16}{MPa} or to untested fluids is unverified.
Second, surface descriptors are sparse in the flow-boiling literature---most
records lack measured roughness and contact angle---so the transferred surface
encoder is exercised less here than in pool boiling. Third, the uncertainty
quantification captures epistemic but not aleatoric uncertainty, under-covering
by $\sim\!4.7\times$; calibrated intervals presently require the interim
rescaling, pending a heteroscedastic head. Fourth, a narrow
intermediate-pressure, low-mass-flux pocket (a single source's
$3.5$--\SI{5}{MPa} runs) retains a residual over-prediction of $\sim\!\SI{1.7}{K}$
that reflects a data gap rather than a model defect.
